% ============================================================
%  Informe Quincenal de Avance n.o 2 - TFG AutoPipe
%  Autor: Ricardo Antonio Borbon Mena
%
%  Contenido segun el Apendice 1 de los Lineamientos para los
%  Trabajos Finales de Graduacion V4.1 (CE-5601).
%
%  Compilar con:  pdflatex "Avance Quincenal 2 - Ricardo Borbón.tex"
% ============================================================

\documentclass[11pt,letterpaper]{article}

\usepackage[utf8]{inputenc}
\usepackage[T1]{fontenc}
\usepackage[spanish,es-tabla]{babel}
\usepackage{lmodern}
\usepackage[letterpaper,top=2.4cm,bottom=2.4cm,left=2.6cm,right=2.4cm]{geometry}
\usepackage{xcolor}
\usepackage{colortbl}
\usepackage{booktabs}
\usepackage{array}
\usepackage{tabularx}
\usepackage{enumitem}
\usepackage{fancyhdr}
\usepackage{titlesec}
\usepackage{tikz}
\usetikzlibrary{shapes.geometric}

\definecolor{tecnavy}{HTML}{003366}
\definecolor{hecho}{HTML}{003366}
\definecolor{quincena}{HTML}{C87F0A}
\definecolor{pendiente}{HTML}{D9D9D9}
\definecolor{bordepend}{HTML}{9A9A9A}
\definecolor{gris}{HTML}{666666}
\definecolor{tablehead}{HTML}{E6EEF5}

\newcolumntype{Y}{>{\raggedright\arraybackslash}X}
\newcolumntype{L}[1]{>{\raggedright\arraybackslash}p{#1}}

\setlength{\parindent}{0pt}
\setlength{\parskip}{5pt}
\renewcommand{\arraystretch}{1.25}
\setlist[itemize]{leftmargin=1.3em,topsep=3pt,itemsep=3pt}
\setlist[enumerate]{leftmargin=1.5em,topsep=3pt,itemsep=3pt}

\titleformat{\section}{\large\bfseries\sffamily\color{tecnavy}}{\thesection.}{0.6em}{}
\titleformat{\subsection}{\normalsize\bfseries\sffamily}{}{0em}{}
\titlespacing*{\section}{0pt}{16pt}{8pt}
\titlespacing*{\subsection}{0pt}{12pt}{4pt}

\pagestyle{fancy}
\fancyhf{}
\renewcommand{\headrulewidth}{0.4pt}
\fancyhead[L]{\small AutoPipe --- TFG}
\fancyhead[R]{\small Informe de avance n.\textsuperscript{o}~2}
\fancyfoot[C]{\thepage}

\begin{document}

% =========== ENCABEZADO ====================================
\vspace*{0.12\textheight}
\begin{center}
{\LARGE\bfseries\color{tecnavy} Informe Quincenal de Avance}\\[4pt]
{\large Trabajo Final de Graduación}\\[10pt]
\end{center}

\vspace{10pt}
\begin{center}
\begin{tabularx}{0.92\textwidth}{L{3.6cm} Y}
\toprule
\rowcolor{tablehead} \multicolumn{2}{l}{\textbf{Identificación}} \\
\midrule
Estudiante & Ricardo Antonio Borbón Mena \\
Organización & Instituto Tecnológico de Costa Rica, Escuela de Ingeniería
en Computadores \\
Proyecto & Generalización de framework de selección automática de pipelines de
preprocesamiento a múltiples tareas de aprendizaje automático mediante
meta-aprendizaje \\
Período del informe & 31 de agosto al 14 de septiembre de 2026 \\
Informe n.\textsuperscript{o} & 2 \\
Profesor asesor & Dr.Eng. Jeferson González Gómez \\
Profesor tutor & Luis Alberto Chavarría Zamora \\
Fecha de entrega & 14 de septiembre de 2026 \\
\bottomrule
\end{tabularx}
\end{center}

\vspace{20pt}
\begin{center}
\begin{minipage}{0.86\textwidth}
\centering
{\small\color{gris} Este informe retoma la periodicidad quincenal anunciada en el
anterior. En la quincena se cerraron las semanas 7, 8 y 9 del plan de trabajo y se
inició la 10, de modo que el primer experimento ya está ejecutado.}
\end{minipage}
\end{center}

% =========== 1. TAREAS REALIZADAS ===========================
\newpage
\section{Tareas realizadas en el período}

Las tareas se agrupan por el objetivo específico del anteproyecto. Los objetivos 1 y 2
quedaron alcanzados en el período anterior.

\subsection{Objetivo específico 3 --- Recomendador meta-aprendido}

\begin{itemize}
  \item Se analizó de forma exhaustiva la diversidad de los pipelines ganadores,
        evaluando los 144 pipelines del espacio sobre los 25 conjuntos del corpus. El
        óptimo casi nunca es único ---la mediana es de nueve configuraciones empatadas---
        y la meseta es ancha: una mediana de 36 de las 144 configuraciones alcanza el
        95\,\% del óptimo.
  \item Se implementó una regla de desempate canónico entre configuraciones empatadas y
        se regeneró la base de meta-conocimiento: 21 de las 25 entradas cambiaron de
        configuración sin que cambiara ninguna métrica, y los ganadores que omitían la
        codificación, que no se pueden ejecutar en conjuntos con variables categóricas,
        pasaron de siete a cero.
  \item Se implementó el recomendador por vecinos más cercanos. Compara los conjuntos por
        la distancia L1 entre sus meta-características, con un escalado robusto ajustado
        solo con los conjuntos de cada pliegue y la consulta recortada al rango
        observado, y devuelve cinco configuraciones distintas junto con los vecinos que
        aportó cada una, para que la recomendación sea auditable. Una ejecución en seco
        sobre los 25 pliegues confirmó que ninguno ve su propio conjunto entre los
        vecinos.
  \item Se integró el arranque en caliente en el motor de búsqueda bayesiana: las
        configuraciones se validan al construir el motor y se inyectan como primeros
        ensayos. La validación responde a dos modos de fallo reproducidos contra la
        versión actual de Optuna: un índice fuera de rango aborta el estudio completo y
        una clave desconocida se acepta en silencio.
  \item Se amplió el recomendador a cuatro estrategias con el mismo mecanismo: la
        primaria por similitud, una diversa y dos controles, uno por frecuencia que
        ignora las meta-características y una ablación aleatoria. El control por
        frecuencia no estaba en la hoja de ruta y se añadió porque es el que aísla la
        contribución del meta-aprendizaje.
  \item Se evaluaron las cuatro estrategias sobre los 27 candidatos que no entraron al
        corpus, con la rejilla exhaustiva de cada uno. Todas producen configuraciones
        válidas y cercanas al óptimo, y quedan empatadas: se interpretó como un efecto
        piso de la medición y no como un resultado sobre las estrategias.
  \item Se cerró la interfaz pública del recomendador y se agregó un ejemplo ejecutable
        del flujo completo con arranque en caliente, en 19 instrucciones.
\end{itemize}

Estado: alcanzado. El recomendador produce configuraciones válidas y ejecutables para
conjuntos de datos no vistos y está integrado en el motor de búsqueda como ensayos
iniciales, que es lo que el objetivo declara.

\subsection{Objetivo específico 4 --- Validación experimental}

\begin{itemize}
  \item Se congeló el pre-registro del experimento en su versión 1.0 y se etiquetó el
        mismo día que su commit, antes de ejecutar nada. La versión cierra dos
        decisiones ---cinco semillas por condición y la región de pipelines inevaluables
        declarada como covariable en lugar de corregida--- y declara de antemano la
        anchura de la meseta como segunda covariable.
  \item Se implementó el arnés de validación por exclusión de conjunto de datos contra la
        especificación congelada, con 24 pruebas sobre un corpus sintético que recorren
        el protocolo completo sin producir ningún número del experimento.
  \item Se ejecutó el experimento completo el 13 de septiembre: 25 pliegues, cuatro
        condiciones, cinco semillas, 500 corridas y 15\,000 ensayos.
        \textbf{La hipótesis no se confirmó.} La mediana de ensayos hasta el objetivo
        bajó un 33,3\,\% ---el criterio exige 30\,\%---, pero la prueba de Wilcoxon dio
        $p = 0{,}3115$, y esa reducción equivale a pasar de 3 a 2 ensayos. El control por
        frecuencia sí cumple el criterio ($p = 0{,}0239$): no confirma la hipótesis, pero
        indica que un buen valor por defecto compite con la similitud entre conjuntos.
  \item Se extendió el generador de reportes con un reporte comparativo del experimento,
        con 41 pruebas. Transcribe los estadísticos sin recalcularlos, deriva el
        veredicto de las cifras y marca qué contraste es el primario, para que el
        resultado del control no se lea como una confirmación.
\end{itemize}

Estado: en curso. El experimento está ejecutado y reportado; faltan el análisis de sus
causas, la verificación de la generalización por tarea, el análisis estratificado, la
reproducibilidad de punta a punta y la redacción de resultados.

\subsection{Productos intermedios entregados en el período}

\begin{itemize}
  \item Repositorio de software con la librería: 637 pruebas automatizadas ---448 al
        cierre del informe anterior---, 99\,\% de cobertura, y verificación de estilo y
        de tipos sin incidencias.
  \item Resultados del primer experimento y su reporte comparativo.
  \item Informe de TFG: cambios realizados en el capítulo 1, según sus comentarios, y capitulo 2 completado en su totalidad (marco teórico, estado del arte y metodología).
\end{itemize}

% =========== 2. TAREAS PLANEADAS ============================
\newpage
\section{Tareas planeadas para el siguiente período}

Corresponden al período del 15 al 28 de septiembre, semanas 10 y 11 del plan. El
propósito de la quincena es dejar alcanzado el objetivo específico 4, de modo que al
cierre del período se finalice completamente el proyecto, incluyendo la documentación y el informe final.

\begin{itemize}
  \item Redactar la discusión de los resultados del experimento: dónde se pierde la
        significancia, qué dicen los controles y qué explica el tamaño del efecto.
  \item Diagnosticar las causas del resultado nulo con la covariable declarada, y
        comprobar si el instrumento de medición ---el objetivo al 95\,\% y la variable de
        ensayos hasta el objetivo--- tiene resolución suficiente para discriminar. Si el
        diagnóstico apunta al instrumento, presentarle una propuesta antes de cambiar
        nada.
  \item Verificar la generalización por tarea con un caso completo de detección de
        anomalías en imágenes.
  \item Ejecutar el análisis estratificado por sensibilidad al preprocesamiento declarado
        en el pre-registro.
  \item Verificar la reproducibilidad de punta a punta de los resultados reportados.
  \item Redactar el análisis de resultados del informe.
\end{itemize}

% =========== 3. ANALISIS ====================================
\newpage
\section{Análisis de la situación actual y proyecciones}
\label{sec:analisis}

\subsection{Estado respecto del cronograma}

\begin{figure}[h]
\centering
\begin{tikzpicture}[x=1.18cm, y=1cm, font=\small]

  % --- Encabezado: numero de semana y fecha de inicio ------
  \foreach \s/\f in {1/20 jul, 2/27 jul, 3/3 ago, 4/10 ago, 5/17 ago, 6/24 ago,
                     7/31 ago, 8/7 set, 9/14 set, 10/21 set, 11/28 set, 12/5 oct} {
    \node[anchor=south, font=\footnotesize] at (\s-0.53, 1.52) {\s};
    \node[anchor=south, font=\tiny, text=gris] at (\s-0.53, 1.26) {\f};
  }
  \node[anchor=east, font=\footnotesize, text=gris] at (-0.15, 1.58) {Semana};

  % --- Banda de estado por semana del plan -----------------
  \foreach \s in {1,...,9} {
    \fill[hecho] (\s-1, 0.50) rectangle (\s-0.06, 0.90);
  }
  \foreach \s in {10,...,12} {
    \fill[quincena] (\s-1, 0.50) rectangle (\s-0.06, 0.90);
  }
  \foreach \s in {} {
    \fill[pendiente] (\s-1, 0.50) rectangle (\s-0.06, 0.90);
    \draw[bordepend, line width=0.3pt] (\s-1, 0.50) rectangle (\s-0.06, 0.90);
  }
  \node[anchor=east, font=\footnotesize] at (-0.15, 0.70) {Estado};

  % --- Una fila por objetivo especifico --------------------
  % {etiqueta}/{borde superior}/{semana inicial}/{semana final}/{color}
  \foreach \lab/\yt/\a/\b/\c in {%
      {OE 1}/0.30/1/4/hecho,
      {OE 2}/-0.20/5/7/hecho,
      {OE 3}/-0.70/7/9/hecho,
      {OE 4}/-1.20/9/9/hecho,
      {}/-1.20/10/11/quincena,
      {Cierre}/-1.70/12/12/quincena} {
    \foreach \s in {\a,...,\b} {
      \fill[\c] (\s-1, \yt-0.40) rectangle (\s-0.06, \yt);
      \ifnum\pdfstrcmp{\c}{pendiente}=0
        \draw[bordepend, line width=0.3pt] (\s-1, \yt-0.40) rectangle (\s-0.06, \yt);
      \fi
    }
    \node[anchor=east, font=\footnotesize] at (-0.15, \yt-0.20) {\lab};
  }

  % --- Hitos del plan de trabajo ---------------------------
  \foreach \s in {4,7,9,12} {
    \node[diamond, fill=tecnavy, inner sep=1.7pt] at (\s-0.53, -2.35) {};
  }
  \node[anchor=east, font=\footnotesize] at (-0.15, -2.35) {Hitos};

  % --- Linea de hoy ----------------------------------------
  \draw[tecnavy, dashed, line width=0.7pt] (8, 1.06) -- (8, -2.10);
  \node[anchor=south, font=\tiny, text=tecnavy] at (8, 1.04) {hoy};

  % --- Leyenda ---------------------------------------------
  \fill[hecho] (0, -3.10) rectangle (0.34, -2.86);
  \node[anchor=west, font=\footnotesize] at (0.44, -2.98) {Completado};
  \fill[quincena] (3.0, -3.10) rectangle (3.34, -2.86);
  \node[anchor=west, font=\footnotesize] at (3.44, -2.98) {Próxima quincena};
  \fill[pendiente] (7.2, -3.10) rectangle (7.54, -2.86);
  \draw[bordepend, line width=0.3pt] (7.2, -3.10) rectangle (7.54, -2.86);
  \node[anchor=west, font=\footnotesize] at (7.64, -2.98) {Pendiente};

\end{tikzpicture}
\caption{Estado del plan de trabajo sobre sus doce semanas. La banda y las filas marcan el
avance del plan. Los rombos marcan los hitos previstos: núcleo generalizado (S4), base de meta-conocimiento (S7), recomendador
funcional (S9) y cierre del proyecto (S12).}
\label{fig:cronograma}
\end{figure}

Trabajo realizado (azul): la quincena cerró las semanas 7, 8 y 9 del plan, con los hitos
de la base de meta-conocimiento y del recomendador funcional alcanzados, e inició la 10
con el experimento ya ejecutado. 

El calendario, en cambio, apenas llega al inicio de la
semana 9, de modo que el proyecto va adelantado alrededor de una semana según el calendario. Se
registraron 80 horas en la quincena, para un acumulado de 214 de las 336 del proyecto;
hasta la semana 9, que ya cerró, el plan comprometía 224.

Trabajo próximo (naranja): la quincena que abre este informe corresponde a las semanas 10, 11 y 12 del plan, que cierran la validación experimental y la documentación final.

\subsection{Proyección}

Al cierre de la próxima quincena se espera tener alcanzado el objetivo específico 4: la
discusión y el diagnóstico del experimento, el caso de imágenes, el análisis estratificado,
la reproducibilidad verificada y el análisis de resultados redactado en el informe. Por otra parte, 
también se planea finalizar la documentación técnica, el manual de uso y el informe final, antes
de la entrega para revisión del asesor.

\end{document}
